
\section{Methodology}
\label{sec:method}

Given a scene layout generated by DiffuScene \cite{tang2024diffuscenedenoisingdiffusionmodels}, represented as a set of oriented bounding boxes, our goal is to produce a collision-free 3D scene that satisfies fundamental physical constraints. To this end, we introduce a fine-grained collision resolution pipeline that eliminates object interpenetrations through alternating room-bound containment and Separating Axis Theorem (SAT)-based collision repair. As illustrated in Figure~\ref{fig:pipeline}, DiffuScene first takes a text prompt as input and employs a diffusion model to generate an initial layout of oriented bounding boxes. Corresponding 3D object meshes are then retrieved and positioned, scaled, and rotated according to the predicted layouts. Finally, our Collision Resolver iteratively refines the generated bounding boxes by alternating between room-bound containment and OBB collision repair until a physically plausible, collision-free scene is obtained.



\subsection{Preliminaries}
\label{sec:method:prelim}
\begin{figure*}[t]
    \centering
    \includesvg[width=\textwidth,inkscapelatex=false]{Figure1_ACOMPA.svg}
    \caption{Overview of our post-processing pipeline. A relational scene graph (\emph{Front Of}, \emph{Smaller Than}, \emph{Left}, \emph{Close By}) is passed through a Graph Preprocess Module and then to the DiffuScene baseline, which generates an initial bounding-box layout together with retrieved furniture assets. This raw layout is stochastic and frequently contains inter-object interpenetrations and out-of-bound placements (bottom-left inset, red highlights). Our \emph{Collision Resolver} then iteratively alternates between \emph{Room-Bound Containment} and \emph{OBB Collision Repair} (Algorithm~\ref{alg:collision_resolver}) until convergence, yielding a physically plausible, collision-free layout (bottom-right inset, green highlights) that is used to compose the final 3D scene.}
    \label{fig:pipeline}
\end{figure*}
DiffuScene~\cite{tang2024diffuscenedenoisingdiffusionmodels} formulates indoor scene synthesis as a conditional denoising diffusion process over scene layouts. Starting from Gaussian noise, the diffusion model progressively predicts a clean set of object parameters conditioned on the room type and scene context. After the reverse diffusion process converges, the model outputs a structured scene layout consisting of $M$ oriented bounding boxes, each corresponding to a furniture instance.

Each bounding box is parameterized as
\begin{equation}
\mathbf{b}_i=(l_i,h_i,w_i,x_i,y_i,z_i,\sin\theta_i,\cos\theta_i)\in\mathbb{R}^{8},
\end{equation}
where $(x_i,y_i,z_i)$ denotes the center position $\mathbf{p}_i$, $(l_i,h_i,w_i)$ denotes the object dimensions, and $\theta_i$ is the yaw orientation. We further define an \emph{objectness mask} $\mathbf{o}\in\{0,1\}^{M}$ to distinguish valid furniture instances ($o_i=1$) from auxiliary nodes such as the floor and room boundary.

\subsection{Post-Processing Collision Resolver}
\label{sec:method:resolver}

Since DiffuScene does not enforce any implicit or explicit constraints on the denoised layouts so the generated bounding boxes do show signs of interpenetrations and physical violation. We apply a deterministic post-processing step that eliminates pairwise overlaps through iterative projections. 

While the resolver primarily operates in the 2D floor plane ($xz$), we first apply a fast 1D vertical bounding-box check. If two objects are vertically separated (e.g., a monitor resting on a desk), they are immediately excluded from 2D collision resolution. For vertically overlapping pairs, the resolver combines:
(a) \emph{room-bound containment clamping}, and 
(b) \emph{OBB-SAT collision detection and resolution}, which computes the Minimum Translation Vector (MTV).

\begin{algorithm}[t]
\caption{Post-Processing SAT Collision Resolver}
\label{alg:collision_resolver}                        
\begin{algorithmic}[1]
\REQUIRE Layout boxes $\{\mathbf{b}_i\}_{i=1}^{M}$ with centers $\mathbf{p}_i$; mask $\mathbf{o}$; room bounds; max iterations $K_{\max}=500$; clearance $\varepsilon_{\mathrm{push}}$; bias $\beta$
\ENSURE Collision-free layout $\{\mathbf{b}_i'\}$
\FOR{$k = 1, 2, \ldots, K_{\max}$}
    \STATE $\textit{converged} \gets \textbf{true}$
    \STATE \COMMENT{1. Room-bound containment clamping}
    \FOR{each object $i$ with $o_i = 1$} 
        \STATE Project $\mathbf{p}_i$ inside room bounds   
        \IF{$\mathbf{p}_i$ was moved} \STATE $\textit{converged} \gets \textbf{false}$ \ENDIF
    \ENDFOR                                 
    \STATE \COMMENT{2. Object-object collision resolution}
    \FOR{each pair $(i, j)$ with $i < j$, $o_i = o_j = 1$}
        \IF{vertically separated (1D $y$-axis check)}
            \STATE \textbf{continue} 
        \ENDIF
        \STATE Compute corner sets $\mathcal{C}_i, \mathcal{C}_j$ in $xz$-plane
        \STATE Compute overlap depths $\{\delta_k\}_{k=1}^{4}$ via Eq.~\eqref{eq:sat_overlap}
        \IF{$\exists\, k: \delta_k \leq 0$}
            \STATE \textbf{continue} \COMMENT{no collision}
        \ENDIF
        \STATE $\hat{\mathbf{c}}_{ij} \gets (\mathbf{p}_j - \mathbf{p}_i) / \|\mathbf{p}_j - \mathbf{p}_i\|$
        \STATE Select push axis $\hat{\mathbf{a}}^*$ and depth $\delta^*$ via Eq.~\eqref{eq:axis_score}
        \STATE $\Delta \gets (\delta^* + \varepsilon_{\mathrm{push}}) / 2$
        \STATE $\mathbf{p}_i \gets \mathbf{p}_i - \Delta \cdot \hat{\mathbf{a}}^*$
        \STATE $\mathbf{p}_j \gets \mathbf{p}_j + \Delta \cdot \hat{\mathbf{a}}^*$     
        \STATE $\textit{converged} \gets \textbf{false}$
    \ENDFOR
    \IF{$\textit{converged}$}
        \STATE \textbf{break}
    \ENDIF
\ENDFOR
\RETURN $\{\mathbf{b}_i'\}$
\end{algorithmic}
\end{algorithm}
\subsubsection{OBB Collision Detection via SAT}
\label{sec:method:obb_sat}

We detect collisions between oriented bounding boxes using the Separating Axis Theorem (SAT) \cite{10.1145/237170.237244}. For two OBBs with yaw angles $\theta_i, \theta_j$, we test four candidate separating axes. 

Let $\mathbf{a}_k$ ($k = 1, \ldots, 4$) denote these axes, and let $\mathcal{C}_i, \mathcal{C}_j$ denote the sets of four floor-plane corners of boxes $i$ and $j$. For each axis $\mathbf{a}_k$, we project both sets of corners and compute the \emph{overlap depth}:
\begin{equation}
\begin{split}
\delta_k &= \min\!\big(\max(\mathcal{C}_i \cdot \mathbf{a}_k), \max(\mathcal{C}_j \cdot \mathbf{a}_k)\big) \\
&\quad - \max\!\big(\min(\mathcal{C}_i \cdot \mathbf{a}_k), \min(\mathcal{C}_j \cdot \mathbf{a}_k)\big).
\end{split}
\label{eq:sat_overlap}
\end{equation}
If $\delta_k \leq 0$ for any axis, the boxes are not colliding. Otherwise, all four axes yield positive overlaps.

\subsubsection{Alignment-Biased Axis Selection}
\label{sec:method:axis_selection}

The standard SAT approach selects the axis with the minimum overlap depth $\delta^* = \min_k \delta_k$. For interior layouts, this often produces unintuitive sliding. We introduce an \emph{alignment-biased} axis selection that prefers axes aligned with the center-to-center direction $\mathbf{c}_{ij} = \mathbf{p}_j - \mathbf{p}_i$. 

Each axis $\mathbf{a}_k$ is oriented to point from $i$ toward $j$: $\hat{\mathbf{a}}_k = \mathrm{sign}(\mathbf{a}_k \cdot \mathbf{c}_{ij}) \cdot \mathbf{a}_k$. We compute an alignment score:
\begin{equation}
\alpha_k = \big|\hat{\mathbf{a}}_k \cdot \hat{\mathbf{c}}_{ij}\big|, \qquad
s_k = \frac{\tilde{\delta}_k}{\alpha_k^{\beta} + \epsilon}
\label{eq:axis_score}
\end{equation}
where $\beta=2$ and $\epsilon$ is a small constant. Crucially, $\tilde{\delta}_k$ represents a \emph{boundary-penalized} overlap depth. If displacing an object forces it outside the room boundary, we apply a massive scalar penalty ($1000\times$) to $\delta_k$. Axes with $\alpha_k < 0.15$ are discarded. The axis with the lowest score $s_k$ is selected as the push direction.



Given the selected push axis $\hat{\mathbf{a}}^*$ and depth $\delta^*$ for a colliding pair $(i, j)$, we symmetrically displace both objects by $(\delta^* + \varepsilon_{\mathrm{push}}) / 2$ along $\hat{\mathbf{a}}^*$:
\begin{equation}
\mathbf{p}_i \gets \mathbf{p}_i - \frac{\delta^* + \varepsilon_{\mathrm{push}}}{2}\, \hat{\mathbf{a}}^*, \qquad
\mathbf{p}_j \gets \mathbf{p}_j + \frac{\delta^* + \varepsilon_{\mathrm{push}}}{2}\, \hat{\mathbf{a}}^*
\label{eq:push}
\end{equation}
where $\varepsilon_{\mathrm{push}} = 0.02\,\text{m}$ provides clearance. Algorithm~\ref{alg:collision_resolver} iterates over all pairs until no collisions remain or $K_{\max} = 500$ is reached.
    